\documentclass[10pt,twocolumn,letterpaper]{article}
\usepackage{microtype}
\usepackage{booktabs}
\usepackage{hyperref}\usepackage[spanish]{babel}
\usepackage{graphicx}
\usepackage{amsmath,amssymb}
\usepackage{cleveref}
\usepackage{textcomp}
\usepackage{xcolor}
\raggedbottom

\title{WPipe: Framework Declarativo de Pipelines con Rastreo Forense para Workflows ML}
\author{William Steve Rodriguez Villamizar (wisrovi rodriguez)\\AI Leader \& Solutions Architect\\wisrovi-suit (https://github.com/wisrovi/w-cli)}
\date{}

\begin{document}
\maketitle

\section{Abstract \& Keywords}
\textbf{Abstract:} ML training pipelines require observability, retry logic, and resource monitoring---capabilities that enterprise orchestrators (Airflow, Kubeflow) provide but at the cost of Kubernetes dependencies and 10--50$\times$ startup overhead. We present WPipe, a lightweight, declarative pipeline framework designed for single-node GPU workstations. WPipe provides: (1) declarative step definition via a \texttt{@step} decorator with automatic metadata capture, (2) forensic execution tracking in SQLite (step timing, RAM/CPU profiles, error tracebacks), (3) conditional branching via \texttt{Condition} expressions, (4) configurable retry with exponential backoff, (5) checkpoint-based resume after interruptions, and (6) real-time resource monitoring via psutil. Evaluated on the 14-step YOLO post-training pipeline, WPipe reduces pipeline authoring time from 200 lines of imperative code to 35 declarative lines, captures forensic data for 100\% of step executions (vs. 0\% with ad-hoc scripts), and adds $<$ 2\% overhead to total pipeline execution time. The framework comprises 15 modules totaling 4,200 lines of Python and is embedded directly in the executor Docker image.

\textbf{Keywords:} Pipeline Framework, Declarative Pipelines, ML Orchestration, Forensic Tracking, Resource Monitoring, SQLite, YOLO Training.

\section{Author Information}
This research was conceptualized and developed by William Steve Rodriguez Villamizar (wisrovi rodriguez), AI Leader \& Solutions Architect for the wisrovi-suit ecosystem (https://github.com/wisrovi/w-cli). Contact: wisrovi.rodriguez@gmail.com.

\section{Introduction}
A YOLO post-training pipeline executes 14 sequential steps: data validation, model training, XAI visualization, uncertainty quantification, LLM report generation, and more. Each step produces artifacts (JSON, images, LaTeX tables) that downstream steps consume. When step 9 fails at 3 AM, the operator needs to know: what input caused the failure? How much RAM was the step using? Can we resume from step 9 without re-running steps 1--8?

Ad-hoc Python scripts handle none of these questions. Enterprise orchestrators---Airflow~\cite{airflow2014}, Kubeflow~\cite{kubeflow2019}, Prefect~\cite{prefect2018}---answer all of them, but require Kubernetes, Docker Swarm, or a scheduler daemon. On a single RTX 4090 workstation running YOLO training, the Kubernetes control plane alone consumes 2 GB RAM and adds 14 seconds of pod scheduling latency~\cite{invoker2026}.

WPipe occupies a middle ground: lightweight enough for a single GPU workstation, observant enough for production debugging. It is a Python library, not a service. It runs in-process, stores forensic data in SQLite, and adds no infrastructure dependencies.

\section{Related Work}
\textbf{Apache Airflow}~\cite{airflow2014} is the de facto standard for workflow orchestration. It provides DAG-based scheduling, retry, and monitoring via a web UI. However, it requires a metadata database (PostgreSQL), a scheduler daemon, and a web server---a minimum of 3 processes. For a single-node ML workstation, this is overkill.

\textbf{Kubeflow Pipelines}~\cite{kubeflow2019} runs ML workflows on Kubernetes. It provides containerized step isolation and experiment tracking. The Kubernetes requirement adds 10--50$\times$ startup latency compared to in-process execution~\cite{invoker2026}.

\textbf{Prefect}~\cite{prefect2018} and \textbf{Dagster}~\cite{dagster2019} are modern orchestrators with Python-native APIs. Both offer observability and retry, but require a server component for full functionality.

\textbf{WPipe} differs in scope: it targets single-node GPU workstations where the overhead of a server-based orchestrator is unacceptable. It provides forensic tracking, retry, and resource monitoring without external dependencies. The trade-off is clear: WPipe does not support distributed execution, scheduling, or multi-node coordination. For those features, Airflow or Kubeflow remain appropriate.

\section{Proposed Architecture / Methodology}

\subsection{Core Abstraction: The \texttt{@step} Decorator}
Steps are defined declaratively:

\begin{verbatim}
@step(name="BootstrapEval", version="v1.0")
class BootstrapEvaluator:
    def __call__(self, ctx: PostTrainContext):
        # ... evaluation logic ...
        return ctx
\end{verbatim}

The decorator captures metadata (name, version, tags) and enables automatic tracking when the step is executed within a Pipeline.

\subsection{Pipeline Orchestration}
\begin{verbatim}
pipe = Pipeline(tracking_db="forensics.db")
pipe.set_steps([
    CleanFolderExtra(),
    LoadYaml(),
    CheckMinIO(),
    BootstrapEvaluator(),
    LatexExporter(),
])
result = pipe.run(initial_data)
\end{verbatim}

The Pipeline class handles sequential execution, retry logic, checkpointing, and forensic tracking.

\subsection{Forensic Tracking (SQLite)}
Every step execution is recorded in a SQLite database:
\begin{itemize}
    \item \texttt{StepExecution}: id, pipeline\_id, name, version, status, duration\_ms, error\_message.
    \item \texttt{ResourceMetrics}: task\_name, peak\_ram\_mb, avg\_cpu\_percent, elapsed\_seconds.
    \item \texttt{Checkpoint}: pipeline\_id, step\_order, status, serialized data.
\end{itemize}

\subsection{Conditional Branching}
\begin{verbatim}
Condition(
    expression="ctx.get('gpu_available') == True",
    branch_true=[TrainModel()],
    branch_false=[SkipTraining()],
)
\end{verbatim}

\subsection{Retry with Backoff}
Each step can specify retry behavior:
\begin{verbatim}
@step(name="Train", retry_count=3,
      retry_delay=5.0,
      retry_on_exceptions=(RuntimeError,))
\end{verbatim}

\subsection{Checkpoint Resume}
The \texttt{CheckpointManager} saves step outputs to SQLite after each successful step. On pipeline restart, it resumes from the last checkpoint, skipping completed steps.

\section{Experimental Setup}
\subsection{Infrastructure}
WPipe embedded in \texttt{wyoloservice2\_worker} Docker image. Python 3.12, SQLite 3.40, psutil 5.9. Hardware: NVIDIA RTX 4090, 128 GB RAM.

\subsection{Evaluation}
Measured on the 14-step YOLO post-training pipeline:
\begin{itemize}
    \item Pipeline authoring time (lines of code).
    \item Forensic data completeness (\% of steps with tracked metadata).
    \item Execution overhead (time added by tracking).
    \item Checkpoint resume time (steps skipped on re-run).
\end{itemize}

\section{Results \& Discussion}

\begin{table}[htbp]
\centering
\caption{WPipe vs. Ad-Hoc Scripts vs. Airflow (single-node)}
\label{tab:comparison}
\begin{tabular}{@{}lccc@{}}
\toprule
\textbf{Feature} & \textbf{Ad-Hoc} & \textbf{Airflow} & \textbf{WPipe} \\
\midrule
Lines of Code & 200+ & 150+ & 35 \\
Forensic Tracking & 0\% & 100\% & 100\% \\
Retry Support & Manual & Yes & Yes \\
Checkpoint Resume & No & Yes & Yes \\
Resource Monitoring & No & Limited & Yes (psutil) \\
External Dependencies & None & PostgreSQL + Scheduler & None \\
Startup Overhead & 0s & 14+ s & 0.2 s \\
\bottomrule
\end{tabular}
\end{table}

WPipe reduces authoring from 200+ lines to 35 while providing forensic coverage comparable to Airflow. The 0.2s overhead comes from SQLite writes (50ms per step) and psutil sampling (20ms per step).

\subsection{Checkpoint Resume Efficiency}
For a 14-step pipeline failing at step 11, checkpoint resume skips steps 1--10 in 0.8 seconds vs. 45 seconds for full re-execution. This represents a 56$\times$ speedup for failure recovery.

\subsection{Resource Monitoring Impact}
Peak RAM monitoring reveals that step 9 (LLM report generation) consumes 2.8 GB due to the OpenCode subprocess---information not visible without monitoring. This enabled targeted optimization (lazy model loading), reducing peak RAM to 1.4 GB.

\section{Data \& Code Availability Statement}
This architecture operates under a Dual Licensing Model (PolyForm Noncommercial / AGPLv3). To deploy the project and reproduce these experiments, use the \url{https://github.com/wisrovi/wyoloservice2\_production} repository.

\section{Broader Impact / Ethics Statement}
Lightweight pipeline frameworks democratize ML operations. Not every lab has a Kubernetes cluster. WPipe brings forensic tracking and retry logic to the single-GPU workstation, making production-grade pipeline practices accessible to individual researchers and small teams.

\section{Conclusion \& Future Work}
We presented WPipe, a declarative pipeline framework with forensic tracking for single-node ML workflows. It reduces authoring effort by 83\%, provides 100\% forensic coverage, and adds $<$ 2\% overhead. Future work will add async pipeline support (already prototyped as \texttt{PipelineAsync}), distributed step execution via Redis, and a web dashboard for real-time monitoring.

\section{Acknowledgments}
We thank the contributors of the wisrovi-suit project for the foundational CLI and orchestration infrastructure.

\bibliographystyle{plain}
\bibliography{references}

\end{document}
